\documentclass[11pt,a4paper]{article}

% ---------------------------------------------------------
% PREAMBLE (stable, warning-free, Overleaf-safe)
% ---------------------------------------------------------
\usepackage[utf8]{inputenc}
\usepackage[T1]{fontenc}
\usepackage{lmodern}
\usepackage{amsmath, amssymb, amsfonts}
\usepackage{bm}
\usepackage{geometry}
\usepackage{graphicx}
\usepackage{hyperref}
\usepackage{cite}
\usepackage{authblk}
\usepackage{physics}
\usepackage{mathtools}

\geometry{margin=2.4cm}

\title{\textbf{Companion LXXXVI: Neural Joint Likelihoods for RDCC — Combining Transport, Minkowski Functionals, and Shear Power in a Unified Flow-Based Model}}
\author{Michael Lehmann \\ \texttt{mi.lehmann@gmx.de}}
\date{April 2026}

\begin{document}
\maketitle

\begin{abstract}
Weak-lensing surveys provide multiple complementary probes of the large-scale 
structure, including topological persistence diagrams, Minkowski functionals, and 
the shear power spectrum. Previous Companions developed transport-based 
statistics, multi-probe Fisher forecasts, and cross-correlation analyses for the 
Relaxation-Driven Cyclic Cosmology (RDCC). In this work, we introduce a unified 
neural joint likelihood that combines these observables into a single 
flow-based generative model. The likelihood is constructed using normalizing 
flows trained on multi-probe summary vectors, enabling differentiable inference 
of the relaxation parameter $\alpha$. This framework captures non-Gaussian 
topology, morphological structure, and two-point correlations simultaneously, 
yielding a powerful and flexible tool for RDCC parameter estimation in 
Euclid-like surveys.
\end{abstract}
% ---------------------------------------------------------
\section{Introduction}

Weak-lensing convergence maps contain rich information about the large-scale 
structure. Different observables probe different aspects of this structure:

\begin{itemize}
    \item transport-based persistence-diagram statistics capture non-Gaussian 
          topological information,
    \item Minkowski functionals quantify morphological features,
    \item the shear power spectrum $C_\ell^{\gamma\gamma}$ captures two-point 
          correlations.
\end{itemize}

Previous Companions developed each of these probes individually and constructed 
pairwise joint analyses. The present Companion unifies all three into a single 
neural joint likelihood, enabling coherent and differentiable inference of the 
RDCC relaxation parameter $\alpha$.
% ---------------------------------------------------------
\section{Multi-Probe Summary Vector}

We define the combined summary vector



\[
    \mathbf{s}
    =
    \left\{
        \mathcal{T}_{\ell},
        V_{0}(\nu),
        V_{1}(\nu),
        V_{2}(\nu),
        C_{\ell}^{\gamma\gamma}
    \right\},
\]



where:

\begin{itemize}
    \item $\mathcal{T}_{\ell}$ are transport-based topological statistics,
    \item $V_{i}(\nu)$ are Minkowski functionals,
    \item $C_{\ell}^{\gamma\gamma}$ is the shear power spectrum.
\end{itemize}

This vector captures topology, morphology, and two-point structure simultaneously.
% ---------------------------------------------------------
\section{Flow-Based Likelihood Model}

We construct a normalizing flow $f_{\theta}$ such that



\[
    \mathbf{z} = f_{\theta}(\mathbf{s})
\]



maps the summary vector to a Gaussian latent space. The likelihood is



\[
    \mathcal{L}(\mathbf{s} \mid \alpha)
    =
    \mathcal{N}\!\left(
        f_{\theta}(\mathbf{s});
        \mu(\alpha),
        \Sigma(\alpha)
    \right)
    \left|
        \det \frac{\partial f_{\theta}}{\partial \mathbf{s}}
    \right|.
\]



The flow is trained on simulations across a grid of $\alpha$ values.
% ---------------------------------------------------------
\section{Differentiable RDCC Inference}

Because the likelihood is differentiable with respect to $\alpha$, we can perform:

\begin{itemize}
    \item gradient-based maximum-likelihood estimation,
    \item Hamiltonian Monte Carlo,
    \item variational inference,
    \item neural posterior estimation.
\end{itemize}

The derivative



\[
    \frac{\partial \mathcal{L}}{\partial \alpha}
\]



is computed via automatic differentiation through the flow.
% ---------------------------------------------------------
\section{Advantages of the Unified Approach}

The neural joint likelihood provides:

\begin{itemize}
    \item a coherent combination of topology, morphology, and two-point structure,
    \item improved constraints on $\alpha$ through multi-probe synergy,
    \item a flexible generative model for summary statistics,
    \item robustness to non-Gaussianity and survey noise,
    \item differentiable inference pipelines for Euclid-like surveys.
\end{itemize}

This framework generalizes classical Gaussian likelihoods and captures 
non-linear, non-Gaussian RDCC signatures.
% ---------------------------------------------------------
\section{Conclusion}

We introduced a unified neural joint likelihood that combines topological 
transport statistics, Minkowski functionals, and the shear power spectrum into a 
single flow-based generative model. This framework captures complementary 
information from topology, morphology, and two-point correlations, enabling 
powerful and differentiable inference of the RDCC relaxation parameter $\alpha$. 
The present Companion completes the transition from geometric theory to 
state-of-the-art multi-probe likelihood modeling for Euclid-like surveys.

\bibliographystyle{apsrev4-2}
\nocite{*}
\bibliography{references}


\end{document}
